\section{Why deleting the initializer is a different problem}
\label{sec:scalar-obstruction}

One might try to remove the initializer argument altogether. This section
gives a new scalar-specific obstruction to that route. The argument is about
standard term functions, not about a lower bound for all possible
axiomatizations or proof checkers.

Consider terms generated by variables, $0$, $S$, and the scalar constructor
\[
 F(0;v.h)=S(h(0)),\qquad
 F(Sn;v.h)=h(F(n;v.h)).
\]
Arbitrary nesting and passive numerical parameters are allowed. There is no
primitive numerical equality characteristic or additional arithmetic
operation. The successor may instead be the abbreviation $F(0;v.a)$.

\begin{theorem}[A restriction on every unary scalar term]
\label{thm:unary-bound}
For every unary term $t(n)$ in this language, there is a natural number $C$
such that
\[
 \forall n\in\mathbb N\quad
       \bigl(t(n)\leq C\ \lor\ t(n)\geq n\bigr).
\]
Consequently no such term defines predecessor $\max(n-1,0)$, integer
halving, or any function having unbounded values strictly below its input.
\end{theorem}

The counter-aware tied-seed language in the existing
\texttt{local\_scopes} package already has a finite-control obstruction, but
it does define predecessor. Thus the unary conclusion above is specific to
removing access to the old counter as well as the independent initializer.
We strengthen the existing finite-control invariant; its original complete
proof is included in the archive as
\file{reference/finite_profiles.tex}.

\begin{definition}
Let $\mathcal U$ be the unary functions $f:\mathbb N\to\mathbb N$ for which
there is $C$ satisfying $f(n)\leq C$ or $f(n)\geq n$ at every input.
\end{definition}

\begin{lemma}[Closure properties]
Constants, the identity, and successor belong to $\mathcal U$. The class is
closed under composition and under pointwise selection from a fixed finite
family of its members.
\end{lemma}
\begin{proof}
For finite selection, take the maximum of the finitely many bounds. For
composition $f\circ g$, fix bounds $C_f,C_g$. If $g(n)\leq C_g$, then the
output is bounded by $\max_{j\leq C_g}f(j)$. Otherwise $g(n)\geq n$, and
the output is either at most $C_f$ or at least $g(n)\geq n$.
\end{proof}

\begin{lemma}[The orbit lemma]
If $h\in\mathcal U$, then the sequence
\[
 r(0)=1+h(0),\qquad r(n+1)=h(r(n))
\]
viewed as a unary function of $n$, also belongs to $\mathcal U$.
\end{lemma}
\begin{proof}
A finite orbit is bounded. An infinite orbit has no repeated value, since
a repeat in a deterministic orbit makes its future periodic. Fix a bound
$C$ for $h$. A strictly decreasing transition must end in
$\{0,\ldots,C\}$, so an injective orbit has only finitely many such
transitions. After the last one, the orbit increases strictly. Eventually
it exceeds every value in the earlier finite prefix. For all sufficiently
large $n$, its first $n+1$ distinct values therefore lie in
$[0,r(n)]$, and $r(n)\geq n$. Increase a constant to cover the finitely many
earlier outputs.
\end{proof}

\begin{proof}[Proof of Theorem~\ref{thm:unary-bound}]
Strengthen the finite-control invariant as follows. Every positive-arity
term has one active numerical coordinate; each other coordinate is observed
through a finite unary feature map; once those finite observations are fixed,
the output is a unary profile of the active coordinate. Require every such
profile to belong to $\mathcal U$. Constants and projections satisfy this
property. In a composition, fixing the inactive features leaves a pointwise
finite selection of compositions of old profiles. The closure lemma applies.

For $F$, first treat its counter as an independent input. If the active
coordinate of its body is the state $v$, fixing the finitely many passive
features gives finitely many updates $h_c\in\mathcal U$. Their orbit profiles
belong to $\mathcal U$ by the orbit lemma, and the counter is the new active
coordinate.

If instead the active coordinate of the body is a passive parameter $p$,
write
\[
 h(v,p,\vec q)=H_{\beta(v),\gamma(\vec q)}(p),
\]
where $\beta,\gamma$ take finitely many values and every $H_{b,c}$ belongs
to $\mathcal U$. Form a finite catalogue $V_j(p)$ of all update profiles
$H_{b,c}(p)$ and all initial profiles $1+H_{\beta(0),c}(p)$. Every orbit
state is a catalogue value, and every catalogue member belongs to
$\mathcal U$.

The finite feature
\[
 \kappa(p)=(\beta(V_j(p)))_j
\]
determines every transition of the catalogue indices once
$c=\gamma(\vec q)$ is fixed. For every possible $c,k$, let
$\theta_{c,k}(n)$ be that finite-state process after $n$ updates. Pack all
these finitely many processes into the finite unary feature
$\eta(n)=(\theta_{c,k}(n))_{c,k}$. Then
\[
 F(n;v.h)=V_{\eta(n)_{\gamma(\vec q),\kappa(p)}}(p).
\]
For fixed features of $n,\vec q$, this is a pointwise finite selection of
the catalogue profiles of $p$, so it belongs to $\mathcal U$. This argument
uses only the displayed finite observations, not a numerical equality
oracle. Substitute the actual counter term by the already established
composition closure. Structural induction is complete. At arity one there
are no inactive coordinates, yielding the asserted unary conclusion.
\end{proof}

For predecessor, choose $n>C+1$; its value is then greater than $C$ and less
than $n$, a contradiction. Halving fails at arbitrarily large $n>2C+1$.
The theorem by itself concerns term definitions. It must not be presented
as an obstruction to arbitrary graph interpretations or changed domains.
Those broader questions are treated separately in the existing relational
obstruction, also supplied in the archive.

\subsection{Independence of the original two scalar computations}

Neither computation equation of the original initialized scalar iterator
can simply be deleted, even if full induction is supplied. In the standard
natural numbers, interpret every fixed body compositionally.

If $I_0$ is omitted, put $I(n,a;h)=h^n(0)$. The successor equation holds,
but $I(0,1;u.u)=0\ne1$. If $I_S$ is omitted, put $I(n,a;h)=a$. The zero
equation holds, but the successor equation fails for $h(u)=Su$.
Both structures have ordinary numerical coverage, separation, and every
first-order induction instance. Thus the canonical-$R$ construction succeeds
by changing the representation; it does not derive a deleted scalar law
from the unchanged remaining axioms.

